% !PRINTED 611
\weberpage{611}

%<<CONT
zahlige lineare Substitution mit der Determinante $\pm 1$ ist (\S.\,147), und folglich sind die beiden Formen

\[ (r\omega+s)\,ec\,\varphi_1,\quad e_1c_1\varphi \]

mit einander associirt, und dies hat zur Folge, dass die beiden Formen $\varphi,\varphi_1$ in dem in \S.\,152 definirten Sinne äquivalent sind. Damit haben wir bewiesen:

2. \gesperrt{Wenn die quadratischen Irrationalzahlen $\omega,\omega_1$ äquivalent sind, so sind auch die entsprechenden Functionale $\varphi,\varphi_1$ äquivalent.}

Es lässt sich aber auch der umgekehrte Satz beweisen:

3. \gesperrt{Wenn die Functionale $\varphi,\varphi_1$ äquivalent sind, so sind auch die quadratischen Irrationalzahlen $\omega$ und $\omega_1$ äquivalent.}

Um diesen Beweis zu führen, bezeichnen wir mit $\psi$ ein Functional der Classe, die zu der Classe von $\varphi$ reciprok ist, so dass $\varphi\psi$ mit einer Zahl äquivalent ist. Dann können wir setzen:

\weq{5}{\varphi\psi = \varepsilon\xi,}

worin $\varepsilon$ eine Einheit, $\xi$ eine \gesperrt{Zahl} in $\mathfrak{o}$ ist. Hierin können wir $\varphi$ und $\psi$ beide als Basisformen annehmen und demnach

\weq{6}{\varphi = ec(t_1+t_2\omega),\quad \psi = e'c'(t_1+t_2\eta)}

setzen, worin $e', c', \eta$ dieselbe Bedeutung für $\psi$ haben, wie $e,c,\omega$ für $\varphi$. Wir setzen zur Abkürzung [\S.\,166, (13)]:

\weq{7}{N_a(\varphi) = e^2c = P,\quad N_a(\psi) = e'^2c' = Q.}

Wenn wir mit $\psi'$ die zu $\psi$ conjugirte Form bezeichnen, und für den Augenblick die Einheitsform

\[ c't_1^2 + b't_1t_2 - a't_2^2 = \varepsilon_1 \]

setzen, so folgt aus (5) durch Multiplication mit $\psi'$:

\[ Q\varepsilon_1\varphi = \varepsilon\xi\psi', \]

und folglich

\weq{8}{\frac{Q\varphi}{\xi} = \frac{\varepsilon}{\varepsilon_1}\psi' = \chi,}

worin $\chi$ und $\psi'$ associirt sind. Setzen wir

\weq{9}{\frac{Qec}{\xi} = \beta_1,\quad \frac{Qec\omega}{\xi} = \beta_2,}
%CONT>>
